% ===================================================================
%  TAULA N° 14 — Lo Carrièr. — Los Comerçants.
%  Traduction 2026 d'apres texto/io, controlee sur texto/fr et
%  texto/ca. Consigne : voir texto/oc/10-taula-01.tex.
%
%  LE TABLEAU LE PLUS CHARGE DU LIVRET : quatre-vingt-dix-neuf
%  renvois. Le suedois et le finnois y avaient chacun deux inversions
%  a defaire ; l'occitan n'en a pas une, comme le catalan.
% ===================================================================

%%K t14-apar-1 apar
\VUsaut{4.0mm}
\VUtitre{15.0pt}{60}{TAULA N° 14}
\VUsaut{3.0mm}

%%K t14-tit-1 sub
\VUpk{11.4pt}{40}{Lo Carrièr. --- Los Comerçants.}
\VUsaut{3.0mm}

%%K t14-01-1 p
1. --- Mon amic Joan, que demòra a la campanha, es vengut passar tres jorns amb
ma familha. \VUgras{Fins} ara aviá pas jamai agut l'escasença de veire una granda
vila, de biais que passam totes los jorns a passejar e ieu li explici çò que
coneis pas.

%%K t14-02-1 p
2. --- D'en primièr anèrem veire l'\VUgras{observatòri} \textsuperscript{(1)}, que
l'interessèt fòrça. En tornar pel \VUgras{carrièr} \textsuperscript{(3)} ont son
los \VUgras{banhs turcs} \textsuperscript{(2)}, nos crosèrem amb un
\VUgras{enterrament} \textsuperscript{(4)}. Al cap del \VUgras{seguici funèbre}
caminava lo \VUgras{clergat}, davant lo \VUgras{corbilhard} ont s'èra pausada la
\VUgras{caissa}. Fòrça amics seguissián la familha en \VUgras{dòl}.

%%K t14-02-2 p
Quand lo seguici foguèt passat, crosèrem lo carrièr davant la granda
\VUgras{cerveseriá} \textsuperscript{(5)}, ont fòrça \VUgras{clients} bevián de
\VUgras{cervesa} sus la terrassa, a l'abric de la \VUgras{marquisa}
\textsuperscript{(6)} veirada.

%%K t14-03-1 p
3. --- Joan foguèt intrigat per l'\VUgras{escut d'armas} pausat entre la botiga
del \VUgras{comerçant de vins e licors} \textsuperscript{(7)} e aquela del
\VUgras{vendeire de musica} \textsuperscript{(10)}. Me demandèt çò que significava;
li diguèri que senhalava lo \VUgras{consolat} \textsuperscript{(8)} anglés, e li
mostrèri lo \VUgras{cònsol} \textsuperscript{(9)} drech sul balcon.

%%K t14-tit-2 sub
\VUpk{10.2pt}{40}{En agachar las vitrinas.}
\VUsaut{3.0mm}

%%K t14-04-1 p
4. --- Naturalament nos arrestèrem davant l'exposicion d'\VUgras{instruments de
musica}, d'\VUgras{armòniums}, d'\VUgras{òrgues} e de \VUgras{pianòs}; agachèrem
las cobèrtas illustradas de las \VUgras{particions} e de las \VUgras{pèças} per
pianò, e admirèrem la \VUgras{lira} de fusta daurada que serviá d'ensenha.

%%K t14-05-1 p
5. --- Los nívols de polsa qu'auçavan los \VUgras{escobaires}
\textsuperscript{(11)} e l'\VUgras{escobadoira} \textsuperscript{(12)} nos
obliguèron a passar lèu davant los bèles \VUgras{jòcs de mòbles} del
\VUgras{moblièr} \textsuperscript{(13)} e davant l'exposicion d'\VUgras{estufas} e
de \VUgras{lampas} del \VUgras{ferratièr} \textsuperscript{(14)}.

%%K t14-06-1 p
6. --- En aquel endrech, en efièch, nos calguèt daissar lo trepador, qu'èra
embarrassat pels paquets que se descargavan d'un \VUgras{carri de mudança}
\textsuperscript{(16)} per los menar dins una \VUgras{botiga}
\textsuperscript{(15)} que veniá d'èsser logada.

%%K t14-06-2 p
Aquò nos empachèt de veire las bèlas veituras del \VUgras{carrossièr}
\textsuperscript{(17)}, mas pas de nos arrestar per agachar l'\VUgras{embalaire}
\textsuperscript{(19)} que fasiá una caissa per son vesin, lo \VUgras{vendeire de
bicicletas e d'automobilas} \textsuperscript{(18)}. Puèi, coma èra ja pro tard,
tornèrem a l'ostal.

%%K t14-tit-3 sub
\VUpk{10.2pt}{40}{Un passejament per la vila.}
\VUsaut{3.0mm}

%%K t14-07-1 p
7. --- L'endeman ma maire prepausèt que Joan se faguèsse un retrach, e me demandèt
de l'acompanhar en çò del \VUgras{fotograf} \textsuperscript{(20)}. Aprèp lo
dinnar anèrem a l'\VUgras{estudi} de l'artista, ont nos calguèt esperar una bona
estona. Per passar lo temps agachèrem los \VUgras{retrachs}
\textsuperscript{(22)} pendolats a la paret.

%%K t14-07-2 p
Quand nos arribèt lo torn la causa durèt pas gaire, e tanlèu que mon amic aguèt
acabat de pausar davant la \VUgras{camèra} \textsuperscript{(21)}, davalèrem.

%%K t14-08-1 p
8. --- Al primièr estatge vegèrem, per la pòrta dobèrta de la \VUgras{perruquiá}
\textsuperscript{(23)}, l'oficial que talhava los pels a un client.

%%K t14-08-2 p
Ja tornats al carrièr, passèrem davant lo \VUgras{burèu de tabat}
\textsuperscript{(24)}. Joan foguèt tan divertit pel \VUgras{fanal}
\textsuperscript{(25)} roge e per la \VUgras{granda pipa} que serviá
d'\VUgras{ensenha} \textsuperscript{(26)} que trabuquèt contra dos ancians que
charravan en prenent una \VUgras{presa de tabat} \textsuperscript{(27)}.

%%K t14-tit-4 sub
\VUpk{10.2pt}{40}{La pastissariá.}
\VUsaut{3.0mm}

%%K t14-09-1 p
9. --- Los \VUgras{bilhets} e las \VUgras{monedas} de totes los païses expausats
dins la \VUgras{vitrina} del \VUgras{cambiaire} \textsuperscript{(29)} nos
retenguèron un moment, mas coma èra l'ora de la merenda dintrèrem dins la
\VUgras{pastissariá} \textsuperscript{(31)} sens nos arrestar mai.

%%K t14-10-1 p
10. --- Davant totas las \VUgras{lecarditz} qu'i aviá dins la botiga ---
\VUgras{pastisses, pan d'espècias, gaufretas} e \VUgras{doçors} de tota mena ---
podiam tot bèl just causir. A la fin Joan se decidiguèt pels \VUgras{pastisses de
crèma}, e ieu prenguèri pas qu'una pichona \VUgras{tarta de cerièras}, que las
doçors \VUgras{me fan mal de caissal} e ai pas cap d'enveja de visitar tròp sovent
lo \VUgras{dentista} \textsuperscript{(30)}.

%%K t14-tit-5 sub
\VUpk{10.2pt}{40}{L'escritura a maquina.}
\VUsaut{3.0mm}

%%K t14-11-1 p
11. --- Per mostrar a mon amic una \VUgras{maquina d'escriure}, que n'aviá
ausit parlar mas que l'aviá pas jamai vista, dintrèrem dins lo \VUgras{burèu}
\textsuperscript{(32)} de mon \VUgras{cosin} \textsuperscript{(34)}. Lo trobèrem a
dictar sas letras a sa \VUgras{secretària} \textsuperscript{(35)}, qu'es
\VUgras{estenografa}. A pena acabada una letra, un \VUgras{mecanograf}
\textsuperscript{(33)} la copiava a maquina. Coma voliam pas interrompre lo
trabalh de mon parent, i demorèrem pas que qualques minutas.

%%K t14-12-1 p
12. --- Jos lo burèu de mon cosin demòra un grand \VUgras{relotgièr e joielièr}
\textsuperscript{(37)} qu'es en mai argentièr. Son esplendida botiga conten fòrça
\VUgras{relòtges de paret, pendulas, relòtges de pòcha}, \VUgras{cadenas} e
\VUgras{jòias} preciosas, coma de \VUgras{colars de pèrlas}, de
\VUgras{braçalets} d'aur, de \VUgras{pendents} e de \VUgras{anèls} amb de
\VUgras{pèiras preciosas} e de \VUgras{diamants}. Una de las vitrinas es
consacrada tota a l'\VUgras{argentariá} e conten una granda colleccion de
\VUgras{cobèrts} d'argent.

%%K t14-13-1 p
13. --- En virar lo canton del carrièr, notèri qu'avián cambiat la \VUgras{placa}
\textsuperscript{(36)} qu'es contra lo \VUgras{numèro} de l'ostal. Anavi o dire a
votz nauta quand s'ausiguèt lo son alègre d'una musica militara. Nos metèrem a
córrer cap al regiment, sens pensar que passariá davant las botigas del
\VUgras{capelièr} \textsuperscript{(39)} e del \VUgras{gantièr}
\textsuperscript{(40)}, e que seriam estats tan plan plaçats per lo veire desfilar
en demorant davant la vitrina de l'\VUgras{opticaire} \textsuperscript{(38)},
plena de \VUgras{lunetas de pincèl, lunetas} e \VUgras{binòcles}.

%%K t14-tit-6 sub
\VUpk{10.2pt}{40}{La desfilada.}
\VUsaut{3.0mm}

%%K t14-14-1 p
14. --- Al cap del \VUgras{regiment} \textsuperscript{(41)} marchava lo
\VUgras{tambor major} \textsuperscript{(42)}, seguit dels \VUgras{tambors}
\textsuperscript{(43)}, dels \VUgras{clarons} \textsuperscript{(44)} e de tota la
\VUgras{musica}. Lo \VUgras{general} \textsuperscript{(45)}, qu'èra sus la plaça
amb son ajudant de camp, s'èra arrestat contra lo \VUgras{gisclador}
\textsuperscript{(49)} de la \VUgras{font} \textsuperscript{(48)} per veire la
\VUgras{desfilada}.

%%K t14-14-2 p
\VUgras{Los oficièrs de l'estat major} \textsuperscript{(46)}, lo \VUgras{colonèl}
e lo \VUgras{tenent colonèl} lo saludavan amb l'espasa. \VUgras{Los comandants}
anavan al cap de lors \VUgras{batalhons} \textsuperscript{(47)} e cada
\VUgras{capitani} comandava sa \VUgras{companhiá}, mentre que los
\VUgras{tenents}, los \VUgras{sostenents} e los \VUgras{sosoficièrs} ocupavan lors
plaças costumièras contra lors òmes.

%%K t14-15-1 p
15. --- Fòrça passants s'èran acampats a la \VUgras{crosada}
\textsuperscript{(50)}; los comerçants --- lo \VUgras{parapluèjaire}
\textsuperscript{(51)}, lo \VUgras{tenchurièr} \textsuperscript{(53)},
l'\VUgras{armurièr} \textsuperscript{(54)} --- sortiguèron veire los
\VUgras{soldats}. Totes los vehiculs --- \VUgras{veituras de plaça}
\textsuperscript{(52)}, \VUgras{veituras particularas} \textsuperscript{(57)} e
tanben la \VUgras{cuba d'arrosatge} \textsuperscript{(56)} --- s'arrimbèron al
trepador.

%%K t14-tit-7 sub
\VUpk{10.2pt}{40}{Scènas del carrièr.}
\VUsaut{3.0mm}

%%K t14-15-2 p
Un \VUgras{viatjaire de comèrci} \textsuperscript{(55)} que portava a la man sas
\VUgras{valisetas de mòstras} crosèt lo carrièr a bon pas, e las personas
assetadas sus l'\VUgras{imperiala} \textsuperscript{(59)} de l'\VUgras{omnibús}
\textsuperscript{(58)} se virèron per gausir de l'espectacle. Un paure
\VUgras{cantaire de carrièra} \textsuperscript{(60)} amb son \VUgras{orguenet}
\textsuperscript{(61)} deguèt interrompre sa \VUgras{cançon}, que lo bruch dels
tambors negava sa votz.

%%K t14-16-1 p
16. --- Quand lo regiment arribèt davant la \VUgras{botiga de teissuts}
\textsuperscript{(64)}, las \VUgras{cordurièiras} \textsuperscript{(63)} e los
\VUgras{sartres} \textsuperscript{(62)} daissèron lors \VUgras{maquinas de
cordurar} e lor labor per s'abocar a la fenèstra. \VUgras{Lo comerçant}
\textsuperscript{(65)}, que veniá de pausar de \VUgras{vestits}
\textsuperscript{(66)} novèls dins sa \VUgras{vitrina}, deguèt retirar las pèças
de \VUgras{drap} \textsuperscript{(67)} e de \VUgras{tela} expausadas sul
trepador, contra la \VUgras{boca d'escorreguda} \textsuperscript{(68)}, per paur
que lo mond las tombèssan; quitament un vièlh \VUgras{marchand ambulant}
\textsuperscript{(69)}, al qual una portièra marchandava de caucetas, deguèt se
ficar dins un portal per acabar sa venda.

%%K t14-16-2 p
Acompanhèrem lo regiment fins a la casèrna \VUgras{e}, fòrça contents de nòstre
jorn, tornèrem a l'ostal a l'ora de sopar.

%%K t14-tit-8 sub
\VUpk{10.2pt}{40}{Mestièrs e professions divèrses.}
\VUsaut{3.0mm}

%%K t14-17-1 p
17. --- L'endeman mon paire s'ofriguèt de nos menar visitar l'estampariá del
jornal local. O acceptèrem amb entosiasme.

%%K t14-17-2 p
L'\VUgras{estampaire} es tanben \VUgras{librari} \textsuperscript{(70)},
\VUgras{editor}, \VUgras{papetièr} e \VUgras{religaire}. A installat al primièr
estatge la \VUgras{redaccion} \textsuperscript{(71)} del jornal, e tanben
l'\VUgras{obrador de gravadura}, ont trabalhan los \VUgras{litografs}
\textsuperscript{(72)} e los \VUgras{gravaires sus coire} \textsuperscript{(73)},
e enfin la sala de composicion, ont son los \VUgras{compositors}
\textsuperscript{(74)}. La \VUgras{sala de maquinas} \textsuperscript{(75)}
pròpriament dicha, las \VUgras{premsas} e las divèrsas maquinas, son al
ras-de-caussada.

%%K t14-tit-9 sub
\VUpk{10.2pt}{40}{Joan fa fòrça questions.}
\VUsaut{3.0mm}

%%K t14-18-1 p
18. --- En sortir de l'estampariá, Joan s'arrestèt, fòrça intrigat, davant una
caisseta de veire montada sus un pal davant la \VUgras{merçariá}
\textsuperscript{(77)}, e demandèt de qué serviá. Mon paire li expliquèt qu'èra un
\VUgras{avertidor d'incendi} \textsuperscript{(76)}. En efièch, tot dins la vila
èra una meravilha per mon amic, de la simpla \VUgras{font de beure}
\textsuperscript{(78)} e del leugièr \VUgras{tricicle de liurason}
\textsuperscript{(80)} que montava un \VUgras{grom} \textsuperscript{(79)} de
liurèa, fins al \VUgras{furgon de pòsta} \textsuperscript{(81)}.

%%K t14-19-1 p
19. --- Mentre que mon paire nos daissava un moment per parlar amb lo
\VUgras{percebeire de contribucions} \textsuperscript{(83)}, que s'èra arrestat a
charrar amb un \VUgras{avocat} \textsuperscript{(82)} e un \VUgras{notari}
\textsuperscript{(84)}, nos divertiguèrem a seguir lo movement del carrièr.
Passava davant nosautres lo mond mai variat: un \VUgras{aprentís de pastissièr}
\textsuperscript{(85)}, que portava dins un panièr un esplendid \VUgras{pastís},
veniá darrièr un \VUgras{pelharòt} \textsuperscript{(86)}; un
\VUgras{alucalums} \textsuperscript{(87)} parlava amb un \VUgras{gasista}
\textsuperscript{(88)}; una \VUgras{mongeta} \textsuperscript{(89)} que menava per
la man una orfanèla tornava a son convent.

%%K t14-tit-10 sub
\VUpk{10.2pt}{40}{De retorn a l'ostal.}
\VUsaut{3.0mm}

%%K t14-20-1 p
20. --- Just quand mon paire tornava amb nosautres, Joan daissèt anar un
esclafiment de rire en veire las caras mascaradas e l'estranha fatcha de dos
\VUgras{escobachiminèas} \textsuperscript{(91)} negres de sutja.

%%K t14-21-1 p
21. --- Ieu voliái crompar una planta per ma maire a la \VUgras{florista}
\textsuperscript{(92)}, mas mon paire observèt que pesariá fòrça per la portar e
que valdriá mai crompar d'\VUgras{irangas}. Nos apropchèrem de la
\VUgras{vendeirèla} \textsuperscript{(94)}, e quand la \VUgras{governanta}
\textsuperscript{(93)}, qu'èra a ne causir qualqu'unas per son enfant, aguèt acabat
sa crompa, nos serviguèron.

%%K t14-22-1 p
22. --- De retorn a l'ostal nos rescontrèrem amb mantun conegut: una vièlha amiga
de ma familha, apiejada al braç de sa \VUgras{damaisèla de companhiá}
\textsuperscript{(95)}, l'\VUgras{enginhaire} \textsuperscript{(96)} dels camins e
ponts, e una de mas cosinas amb sa \VUgras{noirissa}
\textsuperscript{(98)}.

%%K t14-22-2 p
Puèi passèrem contra la \VUgras{modista} \textsuperscript{(97)} de ma maire e dos
\VUgras{fraires} \textsuperscript{(99)} estrangièrs, que s'apropchèron per
demandar a mon paire lo camin de l'estacion.
\VUsaut{6.0mm}
\VUfilet{22mm}
